% !PRINTED 692
\weberpage{692}

If, conversely, for a system of functions $\eta_1$, $\eta_2$, \dots, $\eta_n$ the
condition
\weq{10}{S(\eta_r\omega_s) = (r,\ s)}
is satisfied, then $\eta_r = \varepsilon_r$; for if one sets
\[ \eta_r = \sideset{^i}{}\sum x_{i,\,r}\varepsilon_i, \]
it follows from (9) and (10) that $x_{r,\,s} = (r,\ s)$. Accordingly the conditions
(9) and (6) are completely equivalent.

The quantities $\varepsilon_i$ form a basis of $\Omega$, but are not
integral functions. The bases (5) and (2) are called \gesperrt{complementary
to one another}.

The quantities
\weq{11}{\Delta\varepsilon_1,\ \Delta\varepsilon_2,\ \dots,\ \Delta\varepsilon_n}
are integral functions, and if $t_1$, $t_2$, \dots, $t_n$ are functional variables,
then
\weq{12}{\Delta(\varepsilon_1t_1 + \varepsilon_2t_2 + \cdots + \varepsilon_nt_n) = \Delta\varepsilon}
is the greatest common divisor of the functions (11).

If
\weq{13}{r = (z - c_1)(z - c_2)(z - c_3)\dots}
is the product of all the linear factors of $\Delta$ that are distinct from
one another, and
\weq{14}{\varkappa = \pi_1\pi_2\pi_3\dots}
the product of the distinct prime functionals dividing $\Delta$,
then we can put every integral function $\varrho$ of $z$ in $\Omega$ that is
divisible by $r$ into the form:
\[ \varrho = r(x_1\varepsilon_1 + x_2\varepsilon_2 + \cdots + x_n\varepsilon_n), \]
in which $x_1$, $x_2$, \dots, $x_n$ are rational functions of $z$. Then
one obtains by (9)
\[ rx_i = S(\varrho\omega_i), \]
and since $\varrho\omega_i$ is divisible by every prime factor of $r$, by the
theorem of § 184, 7. $S(\varrho\omega_i)$ is an integral function in $Z$
divisible by $r$, and consequently $x_i$ an integral rational function.

From this it follows that
\[ \Delta\varrho = r(x_1\Delta\varepsilon_1 + x_2\Delta\varepsilon_2 + \cdots + x_n\Delta\varepsilon_n) \]
is divisible by $r\varepsilon\Delta$, and that accordingly $r\varepsilon$ is an integral
functional dividing $\varrho$.

If $\pi$ were a prime factor of $r\varepsilon$ not contained in $\varkappa$, then
one could assume $\varrho$ to be such that it is not divisible by $\pi$,
%CONT>>
